% ============================================================
% CAS/Python ενότητες — Κεφάλαιο 10: Αόριστο Ολοκλήρωμα
%
% QR code: qr_c58c73614409.png  (→ latex/qrcodes/)
% URL: https://nikmatz.github.io/ELADIC_GR/code/ch10/
% ============================================================

\casactivbox%
  {Maxima: Αόριστο Ολοκλήρωμα — Βασικά, Αντικατάσταση, Κατά Παράγοντες}%
  {%
    \label{cas:ch10_integral}
    Χρησιμοποιήστε \texttt{integrate(f,x)} για να υπολογίσετε τα
    παρακάτω αόριστα ολοκληρώματα και επαληθεύστε με \texttt{diff}.
    \begin{enumerate}[label=(\alph*),itemsep=2pt]
      \item Βασικά ολοκληρώματα:
        $\int x^4\,dx$, $\;\int\sqrt{x}\,dx$, $\;\int e^x\,dx$,
        $\;\int\frac{1}{x}\,dx$, $\;\int\sin x\,dx$,
        $\;\int\sec^2 x\,dx$.
        Επαληθεύστε κάθε αποτέλεσμα με \texttt{diff(F,x)}.
      \item \textbf{Αντικατάσταση} $u$: υπολογίστε
        $\int 2x\cos(x^2)\,dx$ θέτοντας $u=x^2$,
        $\;\int 6x(3x^2+1)^5\,dx$,
        $\;\int\sin^3 x\cos x\,dx$.
        Χρησιμοποιήστε \texttt{trigsimp} για απλοποίηση.
      \item \textbf{Κατά Παράγοντες} ($\int u\,dv$): υπολογίστε
        $\int xe^x\,dx$, $\;\int x\sin x\,dx$, $\;\int\ln x\,dx$.
        Επαληθεύστε με \texttt{diff} και συγκρίνετε με χειρωνακτική εφαρμογή.
      \item \textbf{Ρητές Συναρτήσεις}: υπολογίστε
        $\int\frac{1}{x^2-1}\,dx$ (μερικά κλάσματα) και
        $\int\frac{2x+1}{x^2+x-2}\,dx$.
      \item Τριγωνομετρικά: $\int\sin^2 x\,dx$, $\;\int\tan x\,dx$.
        Επαληθεύστε χρησιμοποιώντας ταυτότητες.
    \end{enumerate}
    \smallskip
    \textbf{Εντολές Maxima:}
    \texttt{integrate(f,x)}, \texttt{diff(F,x)},
    \texttt{trigsimp()}, \texttt{ratsimp()},
    \texttt{subst(u,expr,f)}
  }%
  {https://nikmatz.github.io/ELADIC_GR/code/ch10/}%
  {qr_c58c73614409.png}

\subsection*{Δραστηριότητα Python}

\begin{pybox}{Python: Αόριστο Ολοκλήρωμα — SymPy \& SciPy}%
             {https://nikmatz.github.io/ELADIC_GR/code/ch10/}%
             {qr_c58c73614409.png}
\label{py:ch10_integral}
\begin{enumerate}[label=(\alph*),itemsep=2pt]
  \item Με \texttt{sympy.integrate(f, x)} υπολογίστε τα βασικά
    ολοκληρώματα (δυνάμεις, τριγωνομετρικά, εκθετικά, λογαριθμικά).
    Δημιουργήστε πίνακα $f(x) \to F(x)$ με επαλήθευση
    $F'(x)=f(x)$ μέσω \texttt{sympy.diff}.
  \item Αντικατάσταση $u$: υπολογίστε $\int 2x\cos(x^2)\,dx$,
    $\int x e^{x^2}\,dx$, $\int\sin^3 x\cos x\,dx$.
    Σχεδιάστε $f$ και $F$ σε κοινό διάγραμμα.
  \item Κατά Παράγοντες: υπολογίστε $\int xe^x\,dx$,
    $\int x\sin x\,dx$, $\int\ln x\,dx$, $\int x\ln x\,dx$.
    Εκτυπώστε το αναλυτικό αποτέλεσμα με \texttt{pprint}.
  \item Αριθμητική επαλήθευση: για κάθε $F(x)$ υπολογίστε
    το $\int_a^b f(x)\,dx$ και συγκρίνετε με
    \texttt{scipy.integrate.quad} (ανεκτή απόκλιση $<10^{-8}$).
  \item Σχεδιάστε 6 διαγράμματα (ένα για κάθε κατηγορία)
    που δείχνουν ταυτόχρονα $f(x)$ και $F(x)$ με \texttt{matplotlib}.
\end{enumerate}
\medskip
\textbf{Βιβλιοθήκες / Εντολές:}
\texttt{sympy.integrate()}, \texttt{sympy.diff()},
\texttt{sympy.lambdify()}, \texttt{sympy.trigsimp()},
\texttt{scipy.integrate.quad()}, \texttt{matplotlib.pyplot}
\end{pybox}
